\subsection{Toy equivalence audits: scattering delay as queueing and black-hole saturation}
\label{app:scattering_bh_toy_equivalence_audits}

\paragraph{Scope and status.}
\AuditTag This appendix records \emph{toy} equivalence audits in the computer-science (queue/record) language.
It does not attempt to derive scattering from the folding core, and it does not claim a theorem-level equivalence
between scattering and black-hole physics.
The purpose is to provide reproducible artifacts that realize the interface narrative
``unsaturated delay'' versus ``budget-saturated delay (horizon-like)'' in explicit finite models.

\paragraph{Failure-point alignment.}
\AuditTag When the scattering viewpoint is used, the relevant minimal failure point is (S1): approximate unitarity on the declared band
(Appendix~\ref{app:minimal_failure_point_templates} and Appendix~\ref{app:scattering_haag_ruelle_lsz_interface}).
In these toys, loss is explicitly enabled/disabled and summarized by an $\eta_{\min}$ gate, and the record algebra is explicitly declared
as the emitted stream.

\subsubsection{Scattering as a delay-driven queue (toy)}
\label{subsec:scattering_delay_queue_toy}

\AuditTag We simulate a two-channel $2\times2$ proxy matrix $S(\omega)$ on a bounded frequency grid and use the trace Wigner--Smith delay
$\tau_{\mathrm{WS}}(\omega)=\mathrm{Tr}\,Q(\omega)$ as a tick-time service cost for a single-server queue.
The load is controlled by an explicit arrival period (one job per $p$ ticks).

\begin{table}[t]
\centering
\scriptsize
\setlength{\tabcolsep}{6pt}
\renewcommand{\arraystretch}{1.15}
\begin{tabular}{c r r r l r r r r r r r}
\toprule
case & $p$ & scale & loss\_amp & S1 & $\eta_{\min}$ & $\overline{\tau}_{\mathrm{WS}}$ & $\overline{\eta}$ & frac\_abs & max\_q & frac\_q$>q_\star$ & comp\_ratio \\
\midrule
A & 8 & 0.600 & 0.000 & ok & 1.000000 & 6.984603 & 1.000000 & 0.000000 & 6 & 0.000000 & 0.069723 \\
B & 4 & 0.400 & 0.000 & ok & 1.000000 & 7.046241 & 1.000000 & 0.000000 & 11 & 0.000000 & 0.053589 \\
C & 1 & 0.600 & 0.000 & ok & 1.000000 & 7.047209 & 1.000000 & 0.000000 & 4730 & 0.988667 & 0.056699 \\
D & 1 & 0.600 & 0.350 & loss & 0.650722 & 7.047209 & 0.914106 & 0.091339 & 4730 & 0.988667 & 0.079793 \\
\bottomrule
\end{tabular}
\caption{Toy scattering queue: delay-driven service cost and saturation diagnostics. Rows are generated by \texttt{scripts/exp\_scattering\_process\_delay\_queue\_sim.py}.}
\label{tab:scattering_process_delay_queue_toy}
\end{table}

\paragraph{Toy scattering-process simulation (delay$\rightarrow$queue).} \AuditTag We simulate a two-channel scattering proxy via a bounded-frequency $2\times 2$ matrix $S(\omega)$ and use the trace Wigner--Smith delay proxy $\tau_{\mathrm{WS}}(\omega)=\mathrm{Tr}\,Q(\omega)$ to set a discrete tick-time service cost for an internal single-server queue. The external record algebra is a single emitted stream in the stable alphabet $X_6$ (plus an explicit loss token \texttt{X} when a lossy window is enabled). Backlog/saturation is summarized by the maximum queue length and the fraction of ticks exceeding a fixed threshold $q_\star$; loss/approximate-unitarity is summarized by $\eta_{\min}$ on the declared band (S1-style gate).
\paragraph{Audit notes.} \AuditTag Parameters are fixed and deterministic (ticks=6000, grid=81, q\_star=50). The load is controlled by an explicit arrival period (one job every $p$ ticks). The scattering carrier $S(\omega)$ is CAP-selected from an explicit bounded family. If the matching dictionary provides eligible cross-unit targets (M2), the carrier is selected by the multi-dataset CAP audit; otherwise it falls back to benchmark-only anchoring (M1). The purpose is to provide an explicit computational interpretation of the delay dictionary, not to claim a theorem-level scattering construction.
\paragraph{Run notes (deterministic).} \AuditTag Case A: arrival_period=8, scale=0.600, loss_amp=0.000, eta_min=1.000000, max_q=1, frac_over_q=0.000000. Case B: arrival_period=4, scale=0.400, loss_amp=0.000, eta_min=1.000000, max_q=1, frac_over_q=0.000000. Case C: arrival_period=1, scale=0.600, loss_amp=0.000, eta_min=1.000000, max_q=3096, frac_over_q=0.984833. Case D: arrival_period=1, scale=0.600, loss_amp=0.350, eta_min=0.650722, max_q=3096, frac_over_q=0.984833.
\paragraph{Carrier selection notes (deterministic).} \AuditTag Case A: carrier note: carrier_source=registry; omega0=1.35, gamma_ref=1.000000, abslog=0.000000, gap=0.000000. Case B: carrier note: carrier_source=registry; omega0=1.35, gamma_ref=1.000000, abslog=0.000000, gap=0.000000. Case C: carrier note: carrier_source=registry; omega0=1.35, gamma_ref=1.000000, abslog=0.000000, gap=0.000000. Case D: carrier note: carrier_source=registry; omega0=1.35, gamma_ref=1.000000, abslog=0.000000, gap=0.000000.%

\subsubsection{Scattering vs black-hole queue equivalence (toy audit)}
\label{subsec:scattering_bh_queue_equivalence_toy}

\AuditTag We compare the delay-driven scattering queue to a black-hole-like internal queue that emits a stable-label record $w_t\in X_6$.
The arrival period $p$ is shared so that ``unsaturated'' and ``saturated'' regimes can be compared side by side.

\begin{table}[t]
\centering
\scriptsize
\setlength{\tabcolsep}{6pt}
\renewcommand{\arraystretch}{1.15}
\resizebox{\textwidth}{!}{%
\begin{tabular}{r r r r r r r r r r r r r r r}
\toprule
$p$ &
\multicolumn{5}{c}{scattering (unitary)} &
\multicolumn{5}{c}{scattering (lossy)} &
\multicolumn{4}{c}{BH queue ($m=6$)} \\
\cmidrule(lr){2-6}\cmidrule(lr){7-11}\cmidrule(lr){12-15}
& $\eta_{\min}$ & $\overline{\tau}_{\mathrm{WS}}$ & max\_q & frac\_q$>q_\star$ & comp\_ratio
& $\eta_{\min}$ & frac\_abs & max\_q & frac\_q$>q_\star$ & comp\_ratio
& max\_q & frac\_q$>q_\star$ & frac\_leak & comp\_ratio \\
\midrule
8 & 1.000000 & 2.610602 & 1 & 0.000000 & 0.023052 & 0.650722 & 0.081333 & 1 & 0.000000 & 0.059668 & 1 & 0.000000 & 0.078667 & 0.155818 \\
4 & 1.000000 & 2.610602 & 1 & 0.000000 & 0.014763 & 0.650722 & 0.089333 & 1 & 0.000000 & 0.051073 & 1 & 0.000000 & 0.090000 & 0.148252 \\
1 & 1.000000 & 2.610602 & 3096 & 0.984833 & 0.010036 & 0.650722 & 0.094008 & 3096 & 0.984833 & 0.042928 & 6000 & 0.991667 & 0.087000 & 0.141385 \\
\bottomrule
\end{tabular}
}%
\caption{Toy equivalence audit: saturation/backlog and record compressibility proxies for scattering versus a black-hole-like queue. Rows are generated by \texttt{scripts/exp\_scattering\_bh\_queue\_equivalence\_audit.py}.}
\label{tab:scattering_bh_queue_equivalence_toy}
\end{table}

\paragraph{Scattering vs black-hole queue equivalence (toy).} \AuditTag We compare (i) a delay-driven scattering queue (two-channel $S(\omega)$ proxy, with Wigner--Smith delay setting a tick-time service cost) to (ii) a black-hole-like internal queue emitting one stable-type record $w_t\in X_m$ per tick. The shared control parameter is the periodic arrival rate (one job per $p$ ticks). Trap-like saturation is audited by the fraction of ticks where the internal queue length exceeds a fixed threshold $q_\star$, and by the maximum observed queue length. External-record structure is summarized by a simple compressibility proxy (zlib compression ratio) on the emitted record stream.
\paragraph{Scope.} \AuditTag This is an interface-level comparison artifact (CS model). It does not assert a theorem-level equivalence between scattering and black-hole physics.
\paragraph{Determinism.} \AuditTag Parameters are fixed (ticks=6000, q\_star=50, m=6, grid=81). BH knobs are CAP-selected from a bounded family (absorb\_batch in [1, 2, 3, 4], leak\_amp in [0.0, 0.05, 0.091, 0.1, 0.2, 0.35]) to match the lossy-scattering audit targets across arrival periods (RMSE=0.268970, gap=0.000092); selected: absorb\_batch=2, leak\_amp=0.091. Loss is reported via an S1-style $\eta_{\min}$ gate and an absorption fraction proxy in the lossy scattering case.%

\subsubsection{Record-algebra Page-surrogate equivalence (toy audit)}
\label{subsec:scattering_bh_page_surrogate_equivalence_toy}

\AuditTag We compare a Page-like surrogate functional $U(t)$ (remaining ambiguity in bits on record prefixes) between:
(i) a black-hole-like stable-label record hiding folding-fiber micro-information, and
(ii) a scattering-like record hiding a fixed hidden-bit budget per event (toy interface knob).
Both use the same base-$|X_6|$ digit-release schedule template.

\begin{table}[t]
\centering
\scriptsize
\setlength{\tabcolsep}{6pt}
\renewcommand{\arraystretch}{1.15}
\begin{tabular}{r r r}
\toprule
prefix fraction & $U_{\mathrm{BH}}$ (bits) & $U_{\mathrm{SC}}$ (bits) \\
\midrule
0.000 & 0.000000 & 0.000000 \\
0.100 & 15.169925 & 190.000000 \\
0.250 & 37.924813 & 411.584963 \\
0.500 & 71.264663 & 265.614710 \\
0.750 & 34.924813 & 120.000000 \\
0.900 & 9.169925 & 35.614710 \\
1.000 & 0.000000 & 0.000000 \\
\bottomrule
\end{tabular}
\caption{Toy Page-surrogate equivalence audit: remaining ambiguity $U(t)$ on record prefixes. Rows are generated by \texttt{scripts/exp\_scattering\_bh\_page\_surrogate\_equivalence.py}.}
\label{tab:scattering_bh_page_surrogate_equivalence_toy}
\end{table}

\paragraph{Page-surrogate equivalence (toy): scattering vs black-hole record algebra.} \AuditTag We compare a black-hole-like stable-label record (coarse $w\in X_m$ hiding folding-fiber micro-information) to a scattering-like record hiding a fixed number of hidden bits per event. Both are evaluated under the same auditable schedule template (finite-family selection of a mixing schedule) and the same surrogate functional $U(t)$: remaining ambiguity (bits) on record prefixes under progressive release of base-$|X_m|$ digits.
\paragraph{Parameters.} \AuditTag BH anchor $m=6$, $L=48$, $|X_m|=21$, $t_{\mathrm{BH}}=1$, schedule=(scramble\_delay=36, exponent=3). Scattering toy hidden\_bits=10 (g=1024), $t_{\mathrm{SC}}=3$, schedule=(scramble\_delay=36, exponent=3).
\paragraph{Scope.} \AuditTag This is a record-algebra surrogate comparison and does not assert a physical Page-curve derivation.%

